\documentclass[11pt]{article}
\usepackage{geometry}
\usepackage{amsmath}
\usepackage{amssymb}
\usepackage{graphicx}
\usepackage{hyperref}
\usepackage{array}
\usepackage{booktabs}

\geometry{left=2cm,right=2cm,top=2cm,bottom=2cm}

\title{\textbf{XGBOOST: Explicativo para Riesgo Crediticio}}
\author{Sistema de Credito Inteligente}
\date{Mayo 2026}

\begin{document}

\maketitle

\section*{Tabla de Contenidos}
\begin{itemize}
\item Introduccion: Que es XGBoost?
\item Conceptos Fundamentales
\item Data Flow Completo
\item Normalizacion
\item Ingenieria de Caracteristicas
\item Entrenamiento
\item Predicciones
\item Metricas
\item Umbral Optimo
\item Ventajas
\item Glosario
\end{itemize}

\section*{1. INTRODUCCION: QUE ES XGBOOST?}

XGBoost es un algoritmo de machine learning que combina multiples ``expertos debiles'' para crear una ``prediccion fuerte''.

\textbf{Analogia:} Imagina 200 analistas de credito, cada uno con una estrategia simple. Si los 200 votan, su decision conjunta es mucho mas confiable que cualquier experto individual.

XGBoost utiliza:
\begin{itemize}
\item \textbf{Arboles de Decision:} Estructuras simples que hacen preguntas (Si X > Y, entonces...)
\item \textbf{Boosting:} Combina 200 arboles pequenos iterativamente
\item \textbf{Aprendizaje de Errores:} Cada nuevo arbol intenta corregir los errores del anterior
\end{itemize}

\subsection*{Resultado}
\begin{itemize}
\item Velocidad: $\sim$10 milisegundos por prediccion
\item Precision: AUC-ROC 89.83\%
\item F1-Score: 0.73 (Balance precision-recall)
\end{itemize}

\section*{2. CONCEPTOS FUNDAMENTALES}

\subsection*{2.1 Arbol de Decision}

Un arbol de decision es una serie de preguntas:

\begin{verbatim}
  [Inicio: Cliente]
  |
  +-- Es edad > 35?
      |
      +-- Si: Edad_Riesgo = Bajo
      |
      +-- No: Edad_Riesgo = Alto
\end{verbatim}

Cada respuesta (Si/No) divide los clientes en grupos mas homogeneos.

\subsection*{2.2 Boosting}

Boosting significa ``empujar hacia arriba'':

\textbf{Iteracion 1:} Entrena un arbol con todos los datos.

\textbf{Iteracion 2:} Entrena un arbol ENFOCANDOSE en los errores de la iteracion 1.

\textbf{Iteracion 3-200:} Cada arbol intenta corregir los errores acumulados.

\textbf{Prediccion Final:} Suma ponderada de todos los 200 arboles.

\section*{3. DATA FLOW COMPLETO}

\subsection*{Paso 1: Entrada}

Se reciben 8 caracteristicas basicas del cliente:
\begin{enumerate}
\item \texttt{age} - Edad del cliente
\item \texttt{income} - Ingreso anual (en miles USD)
\item \texttt{loan\_amount} - Monto del prestamo (en miles USD)
\item \texttt{interest\_rate} - Tasa de interes (\%)
\item \texttt{employment\_years} - Anos de empleo
\item \texttt{credit\_history\_years} - Anos de historial crediticio
\item \texttt{prior\_default} - Impago previo? (0 o 1)
\item \texttt{debt\_to\_income} - Deuda / Ingreso (ratio)
\end{enumerate}

\subsection*{Paso 2: Ingenieria de Caracteristicas}

A partir de las 8 caracteristicas basicas, se crean 18 nuevas:

\begin{tabular}{|l|l|}
\hline
\textbf{Caracteristica Original} & \textbf{Derivadas (Ejemplos)} \\
\hline
age & age\_square, age\_group\_young, age\_group\_old \\
income & income\_per\_loan, log\_income, income\_stability \\
loan\_amount & loan\_to\_income, loan\_per\_year \\
interest\_rate & rate\_risk\_premium, rate\_volatility \\
\hline
\end{tabular}

\textbf{Por que?} Estos ratios y transformaciones capturan relaciones no lineales que los arboles pueden aprender mejor.

\subsection*{Paso 3: Normalizacion}

StandardScaler transforma cada caracteristica para tener:
\begin{itemize}
\item Media = 0
\item Desviacion estandar = 1
\end{itemize}

\textbf{Formula:}
$$X_{scaled} = \frac{X - \mu}{\sigma}$$

Donde $\mu$ = media, $\sigma$ = desviacion estandar.

\textbf{Beneficio:} XGBoost converge mas rapido cuando todas las caracteristicas estan en la misma escala.

\section*{4. ENTRENAMIENTO}

\subsection*{Configuracion}

\begin{tabular}{|l|r|}
\hline
\textbf{Parametro} & \textbf{Valor} \\
\hline
Numero de Arboles & 200 \\
Profundidad Maxima & 6 \\
Tasa de Aprendizaje & 0.1 \\
Regularizacion L2 & 1.0 \\
Balance de Clases & scale\_pos\_weight \\
\hline
\end{tabular}

\subsection*{Manejo de Desbalance}

En los datos de entrenamiento:
\begin{itemize}
\item 20,186 clientes sin impago (NO-RIESGO)
\item 14,456 clientes con impago (RIESGO)
\end{itemize}

Razon: $\approx$41\% de impagos.

XGBoost ajusta automaticamente los pesos para dar mas importancia a los casos raros (impagos).

\section*{5. PREDICCIONES: EJEMPLO NUMERICO}

\subsection*{Cliente 1: Bajo Riesgo}

\textbf{Datos Basicos:}
\begin{itemize}
\item Edad: 35 anos
\item Ingreso: USD 50,000/ano
\item Prestamo: USD 10,000
\item Tasa: 8.5\%
\item Anos empleo: 5
\item Anos historia crediticia: 10
\item Impago previo?: No (0)
\item Deuda/Ingreso: 0.20
\end{itemize}

\textbf{Paso 1: Ingenieria}
\begin{itemize}
\item age\_square = 1225
\item income\_per\_loan = 5.0
\item loan\_to\_income = 0.20
\item rate\_risk = 8.5 - mean\_rate = -1.2
\end{itemize}

\textbf{Paso 2: Normalizacion}
\begin{itemize}
\item age\_norm = 0.15
\item income\_norm = 0.32
\item loan\_norm = -0.45
\item debt\_to\_income\_norm = -0.33
\end{itemize}

\textbf{Paso 3: XGBoost Prediccion}

200 arboles predicen: ``Riesgo = 0.15'' (15\% probabilidad)

\textbf{Resultado Final:}
\begin{itemize}
\item Probabilidad de Impago: 15\%
\item Recomendacion: APROBAR
\item Risk Score: 15/100
\end{itemize}

\subsection*{Cliente 3: Alto Riesgo}

\textbf{Datos Basicos:}
\begin{itemize}
\item Edad: 45 anos
\item Ingreso: USD 25,000/ano
\item Prestamo: USD 20,000
\item Tasa: 15\%
\item Anos empleo: 0.5
\item Anos historia crediticia: 1
\item Impago previo?: Si (1)
\item Deuda/Ingreso: 0.80
\end{itemize}

\textbf{Normalizacion}
\begin{itemize}
\item debt\_to\_income\_norm = 2.15 (MUY ALTO)
\item employment\_years\_norm = -1.8 (BAJO)
\item prior\_default = 1 (RIESGO)
\item interest\_rate\_norm = 1.95 (ALTO)
\end{itemize}

\textbf{XGBoost Prediccion}

200 arboles predicen: ``Riesgo = 0.78'' (78\% probabilidad)

\textbf{Resultado Final:}
\begin{itemize}
\item Probabilidad de Impago: 78\%
\item Recomendacion: RECHAZAR
\item Risk Score: 78/100
\end{itemize}

\section*{6. METRICAS DE DESEMPENO}

\subsection*{En Dataset de Prueba (6,517 clientes)}

\begin{tabular}{|l|r|}
\hline
\textbf{Metrica} & \textbf{Valor} \\
\hline
AUC-ROC & 0.8983 (89.83\%) \\
Precision & 0.824 (82.4\%) \\
Recall & 0.652 (65.2\%) \\
F1-Score & 0.73 \\
Especificidad & 0.9512 \\
\hline
\end{tabular}

\textbf{Interpretacion:}
\begin{itemize}
\item \textbf{AUC 89.83\%:} El modelo discrimina muy bien entre riesgo y no-riesgo
\item \textbf{Precision 82.4\%:} De los que predecimos como RIESGO, 82.4\% realmente son riesgo
\item \textbf{Recall 65.2\%:} De todos los clientes reales de RIESGO, identificamos 65.2\%
\item \textbf{F1 0.73:} Balance bueno entre precision y recall
\end{itemize}

\section*{7. UMBRAL OPTIMO: 11.6\%}

XGBoost predice probabilidades continuas (0-100\%). Debemos elegir un umbral:

\textbf{Si probabilidad > umbral: RECHAZAR}
\textbf{Si probabilidad <= umbral: APROBAR}

\subsection*{Analisis de Costo-Beneficio}

\begin{tabular}{|l|r|r|}
\hline
\textbf{Umbral} & \textbf{Costo Impago} & \textbf{Costo Falsa Alarma} \\
\hline
5\% & USD 5,000 & USD 50,000 \\
10\% & USD 15,000 & USD 20,000 \\
11.6\% (OPTIMO) & USD 12,000 & USD 12,000 \\
15\% & USD 20,000 & USD 5,000 \\
\hline
\end{tabular}

En 11.6\%:
\begin{itemize}
\item Costo de impago no detectado: USD 12,000 (cliente se fue a la quiebra sin nuestra deteccion)
\item Costo de falsa alarma: USD 12,000 (cliente rechazado que hubiera pagado)
\item Costo total minimizado: USD 24,000
\end{itemize}

\section*{8. VENTAJAS DE XGBOOST}

\begin{enumerate}
\item \textbf{Velocidad:} ~10 ms por prediccion (RAPIDO)
\item \textbf{Precision:} 89.83\% AUC (MUY BUENO)
\item \textbf{Robustez:} Maneja automaticamente desbalance de clases
\item \textbf{Interpretabilidad:} Podemos entender importancia de cada caracteristica
\item \textbf{Escalabilidad:} Maneja millones de clientes
\item \textbf{Estabilidad:} Resultados predecibles y reproducibles
\end{enumerate}

\section*{9. RESUMEN}

\textbf{XGBoost} es un algoritmo que:

1. Toma 8 caracteristicas basicas del cliente
2. Crea 18 caracteristicas derivadas (ratios, transformaciones)
3. Entrena 200 arboles de decision que votan juntos
4. Predice probabilidad de impago con 89.83\% de precision
5. Permite establecer umbral optimo (11.6\%) para decision APROBAR/RECHAZAR

\textbf{Resultado:} Sistema automatizado, rapido y confiable para analisis de credito.

\section*{10. GLOSARIO}

\begin{itemize}
\item \textbf{Arbol de Decision:} Estructura que hace preguntas Si/No para clasificar
\item \textbf{Boosting:} Combinar multiples modelos debiles en uno fuerte
\item \textbf{AUC-ROC:} Area bajo la curva de precision vs recall
\item \textbf{Precision:} De lo que predecimos positivo, cuanto es realmente positivo?
\item \textbf{Recall:} De lo positivo real, cuanto logramos predecir?
\item \textbf{F1-Score:} Promedio armonico de precision y recall
\item \textbf{Normalizacion:} Transformar datos a misma escala
\item \textbf{StandardScaler:} Metodo de normalizacion (media 0, SD 1)
\item \textbf{Desbalance:} Cuando hay mas casos de una clase que otra
\item \textbf{Umbral:} Punto de corte para decision binaria
\end{itemize}

\end{document}
